\section{Decidability and complexity: the two added reductions}
\label{sec:annex-p14-decidability}

\begin{annexitem}{annex:p14.dc.finite-cpl}{thm:p14-dc-finite-cpl}{Finite classical consequence and equivalence}
\begin{proof}
List all distinct atoms occurring in the finite input. Their number and the
lengths of their names are bounded by the input bit length. A valuation is a
bit for each such atom; bottom-up evaluation of each syntax tree and comparison
of atom identifiers take polynomial time in that encoding. A no certificate
for validity is a valuation falsifying the input formula. A no certificate
for consequence is a valuation satisfying every listed premise and falsifying
the conclusion. A no certificate for equivalence is a valuation on which
the two formulas differ. Hence all three languages are in coNP.

For a finite list \(\Gamma=(G_1,\ldots,G_s)\), the classical truth table
for implication gives
\[
 \Gamma\models A
 \quad\Longleftrightarrow\quad
 \models (G_1\land\cdots\land G_s)\to A,
\]
where an empty conjunction is the fixed true abbreviation \(\topv\).
The truth table for biconditional similarly gives
\[
 v(A)=v(B)\text{ for every valuation }v
 \quad\Longleftrightarrow\quad\models A\leftrightarrow B.
\]
Each input tree is copied only a bounded number of times when the abbreviations
are expanded to the primitive syntax. Atom names retain their original bit
strings. The resulting encoding is polynomial in the combined input length.

The coNP-hard validity language is the stated TAUT problem in
\cref{thm:p02-taut-conp-complete}. Map its input \(A\) to the consequence
instance \((\varnothing,A)\); its yes condition is exactly validity. Map it
also to the equivalence instance \((A,\topv)\); again the yes condition is
exactly validity. These are polynomial-time many-one reductions. Validity
already has the cited lower bound, so the three upper and lower bounds match.
If primitive truth is not available, the fixed formula \(q\to q\) expands
\(\topv\), with a fresh atom name of polynomial encoded length. No step
represents an arbitrary infinite premise set by a finite conjunction.
\end{proof}
\end{annexitem}

\begin{annexitem}{annex:p14.dc.horn}{thm:p14-dc-horn}{A logarithmic-space circuit reduction to Horn satisfiability}
\begin{proof}
The upper bound follows from the finite forward-chaining procedure in
\cref{thm:p02-horn-sat}. With a standard explicit clause encoding, parsing,
sorting or comparing finite atom names and carrying out that procedure are
polynomial in bit length. The theorem's sharper literal-reference operation
count is not needed for membership in P.

For hardness use the topologically ordered CVP language from
\cref{def:p04-succinct-problems,thm:p04-succinct-completeness}. Its gates may
include negation, so an unqualified monotone-circuit reduction would not
suffice. Number the \(n\) gates by their positions in the topological list.
For each circuit gate \(i\), introduce two distinct Horn atoms
\(p_i^+\) and \(p_i^-\). They represent, respectively, a true and a false
computed value of that gate. An input constant \(1\) produces the positive
unit \(p_i^+\); a constant \(0\) produces \(p_i^-\). The following are all
positive Horn implications, expanded as Horn clauses:
\[
\begin{array}{c|l}
\text{gate }i&\text{implications emitted}\\ \hline
\neg j&
 p_j^-\to p_i^+,\quad p_j^+\to p_i^-\\
j\land k&
 p_j^+\land p_k^+\to p_i^+,\quad
 p_j^-\to p_i^-,\quad p_k^-\to p_i^-\\
j\lor k&
 p_j^+\to p_i^+,\quad p_k^+\to p_i^+,\quad
 p_j^-\land p_k^-\to p_i^-
\end{array}
\]
For the designated output gate \(o\), emit the additional negative unit
\(\neg p_o^-\). Thus we ask that the false-output rail remain absent, not
that a positive output atom can be guessed arbitrarily.

Let \(H^+\) be the positive clauses just constructed and \(S_*\) their least
forward closure: begin with the constant units and add a clause head whenever
all its antecedents are present. There are only twice as many atoms as circuit
gates, so this reaches a fixed point after finitely many additions. Every
model of \(H^+\) contains \(S_*\), by induction on additions. The valuation
making exactly \(S_*\) true satisfies \(H^+\), because the fixed-point
condition includes the head of every clause whose antecedents it satisfies.

Induction in the circuit's topological order shows
\[
 p_i^+\in S_*\Longleftrightarrow \text{gate }i\text{ has value }1,
 \qquad
 p_i^-\in S_*\Longleftrightarrow \text{gate }i\text{ has value }0.
\]
For constants exactly the correct rail is initially present. For a negation
gate the two implications exchange the predecessor's rails. For conjunction,
the positive rail requires both positive predecessors, whereas the negative
rail requires either negative predecessor. Disjunction is the dual case.
The emitted implications are the only clauses with a given new gate's rails
as heads; thus no later gate can add a spurious rail to an earlier gate.
This proves both directions of the induction, including the fact that the
two rails are never simultaneously present.

If the output is \(1\), then \(p_o^-\notin S_*\), so the least-closure
valuation satisfies the negative unit as well and is a model of the full
Horn formula. If the output is \(0\), the negative rail belongs to \(S_*\)
and therefore to every model of \(H^+\); the negative unit makes the full
formula unsatisfiable. The construction preserves the CVP yes answer exactly.

Gate numbers use \(O(\log n)\) bits. A transducer scans the explicit gate
list, stores the current gate's numbers and type, and emits its constant
number of clauses with indexed rail names. It needs \(O(\log n)\) workspace
and produces polynomially many output bits. Validity of the fixed
topologically ordered encoding is logspace checkable; a malformed input is
mapped to the fixed unsatisfiable Horn conjunction \(q\land\neg q\).
Thus the reduction is total on all encoded words, including no instances
caused by malformed syntax. The cited logspace P-completeness of CVP supplies
P-hardness, and the forward-chaining upper bound completes the classification.
\end{proof}
\end{annexitem}
